\documentclass[10pt]{article}
\usepackage[margin=2.2cm]{geometry}
\usepackage[T1]{fontenc}
\usepackage{lmodern}
\usepackage{textcomp}
\usepackage{amsmath,amssymb}
\usepackage{enumitem}
\usepackage{xcolor}
\usepackage[hidelinks]{hyperref}
\setlength{\parindent}{0pt}
\setlength{\parskip}{4pt}
\newcommand{\rev}[1]{\par\medskip\noindent\textbf{Reviewer comment #1}\par}
\newcommand{\resp}{\par\noindent\textcolor{blue!50!black}{\textbf{Response.}}~}
\newcommand{\chg}{\par\noindent\textcolor{green!35!black}{\textbf{Changes.}}~}

\begin{document}
\sloppy

\begin{center}
{\Large\bfseries Response to the Reviewer --- Second Round}\\[4pt]
Manuscript: \emph{Mission-Driven Corner-Speed Design of a Low-Voltage Switched Reluctance Motor Drive for Electric Bicycles}\\
(previous title: \emph{Corner-Speed-Based Design and Optimization of a Switched Reluctance Motor Drive for Low-Voltage Electric Bicycles})\\[2pt]
Subrahmanyam Tanala
\end{center}

I thank the reviewer for a second careful review and for the clear priority
list in comment~27. Every critical item and every highly recommended item has
been addressed. Prototype measurements are still not available, and the paper
states this throughout. The main additions are:
\begin{enumerate}[nosep]
\item \textbf{Battery voltage.} Battery-voltage sensitivity (30/36/42~V) and a
battery equivalent-circuit model (OCV(SOC), internal resistance, sag) in the
route simulation (Section~VI-E, Table~XI, Eq.~(17)).
\item \textbf{Iron loss.} An iron-loss uncertainty analysis ($\times0.7$ and
$\times1.3$) covering efficiency and route energy (Section~VI-F, Table~XII).
\item \textbf{Coupled-phase check.} Three-phase coupled FE solves at the launch,
200-r/min current-reference, 200-r/min TSF and rated points (Section~IV-D,
Fig.~5).
\item \textbf{Missions.} Three standardized missions (flat urban, hilly urban,
aggressive) (Section~VIII, Table~XVI).
\item \textbf{Corner-speed rule.} A quantitative test over the nine designs
(regression) and over 24/36/48-V redesigns (Section~IX, Fig.~15).
\item \textbf{PM benchmark.} An equal-envelope, equal-boundary analytic
surface-PM benchmark (Section~X-C, Table~XXI). Its result does not favor the
SRM on efficiency, and it is reported as it stands.
\item \textbf{Wording and metadata.} Terminology, novelty claim, ripple
definition, metadata and reproducibility were all corrected.
\end{enumerate}
Table and section numbers refer to the revised manuscript.

%---------------------------------------------------------------------------
\rev{1--4 (strengths; topology wording)}
\resp We thank the reviewer. Following comment~3, Section~III-A now states that
the 12/8 machine ``is selected under the stated envelope, frequency, converter
and performance constraints; it is not the result of a global topology
optimization''.

\rev{5 (``FE validation'' terminology)}
\resp Agreed and corrected throughout. ``Validation'' is now reserved for
comparison with measurements. The FE comparison is called ``FE-based numerical
verification'' or ``comparison with nonlinear FE'' (title of Section~IV,
Section~III-E, the conclusion). The word ``validated'' no longer refers to FE
anywhere in the manuscript.

\rev{6--7 (FE modelling; coupled-phase check for TSF)}
\resp Done. For four operating points we applied the three simulated phase
currents \emph{simultaneously} in the FE model at 24 rotor positions over one
stroke, so that mutual coupling and shared-yoke saturation are included. We
then compared the result with the uncoupled table model used in the drive
simulation:
\begin{itemize}[nosep]
\item launch, current-reference: mean torque $+0.7$\%; ripple 55.7\%
$\rightarrow$ 54.2\%;
\item 200~r/min, 3~N$\cdot$m, current-reference: mean torque $-0.03$\%;
ripple unchanged (194\%);
\item 200~r/min, 3~N$\cdot$m, cubic TSF: mean torque $-0.5$\%; ripple 6.4\%
$\rightarrow$ 7.3\% at the same 24 sampling points;
\item rated, current-reference: mean torque $-1.5$\%; ripple 86.0\%
$\rightarrow$ 87.8\%.
\end{itemize}
The TSF ripple reduction therefore survives coupling. The method is
current-imposed: the uncoupled current waveforms are applied, which is
realistic under hysteresis current control at low speed.
Coupling changes mean torque by at most 1.5\% and ripple by at most 1.8
points.
\chg Section~IV-D (new paragraph, Table~VII, Fig.~5), Section~VII, the
limitations, the abstract and the conclusion.

\rev{8--9 (predictions; launch efficiency)}
\resp Every efficiency, torque, temperature and range figure is now explicitly
labeled as a prediction (abstract, Sections~VI--VIII, conclusion). After the
ripple definition, Section~VI-A adds a paragraph explaining that the 15.9\%
launch efficiency arises because the output is only 54~W while copper loss is
set by the 29-A launch current. It is a short-duration, torque-dominated
condition and not representative of steady-state riding. Its thermal
consequence is quantified in Section~VI-H.

\rev{10 (thermal results are not validated)}
\resp Agreed. The subsection is renamed \emph{Predicted Thermal Envelope} and
opens with: ``The temperatures in this subsection are predictions of an
uncalibrated lumped-parameter model and define a predicted thermal envelope,
not validated temperature rises.''

\rev{11--12 (TSF; torque-ripple metric)}
\resp The ripple metric is now defined mathematically as Eq.~(16):
$\Delta T=(\max T_e-\min T_e)/\bar T_e\times100\%$ over one electrical period.
The definition appears where ripple is first quoted (Section~VI-A), together
with a note that any waveform touching zero once per stroke gives
$\Delta T\geq100$\%.

\rev{13 (single synthetic route)}
\resp Done. Three standardized missions are now evaluated: flat urban (6.4~km),
hilly urban (6.4~km) and aggressive (5.4~km, 6--10\% climbs). Predicted
consumption is 2.5, 5.1 and 10.9~Wh/km (Table~XVI). Section~VIII states that
``a single route figure must not be read as a general consumption prediction''.

\rev{14--15 (battery sag; voltage sensitivity)}
\resp Done in two parts.
\begin{enumerate}[nosep]
\item \textbf{Voltage table.} Table~XI repeats the FE-based analysis at 30, 36
and 42~V, as requested. The launch current (26.5~A) and launch torque do not
depend on voltage. The maximum torque at 25~km/h falls from 4.59 to
2.91~N$\cdot$m at 30~V, which is still about twice the rated torque. The
electrical rated-torque speed limit is 42, 55 and 68~km/h. Rated $\eta_{sys}$
is 77.0, 77.9 and 78.8\%.
\item \textbf{Battery model.} A battery equivalent circuit (Eq.~(17):
representative 10S NMC OCV(SOC) curve, $R_b=0.15~\Omega$) drives
voltage-interpolated FE maps along the missions. Sag changes consumption at the
battery terminals by less than 3\%, and internal battery loss adds 1.5--4\%.
On the aggressive mission at 20\% SOC the loaded voltage reaches 29.8~V, which
touches the controller cut-off.
\end{enumerate}
The estimated range to cut-off with sag is 142, 68 and 32~km for the three
missions.

\rev{16 (battery ripple tolerance claim)}
\resp The unsupported claim is removed. The text now quantifies the internal
heating instead (2.9--6.9~W for the few seconds of a launch). It states that
acceptability over the pack life must be checked against the cell
manufacturer's ripple-current specification.

\rev{17 (no prototype)}
\resp Prototype data are still not available. The paper is framed as an
analytical, FE and system-simulation methodology. The new title says
``Mission-Driven \ldots Design'', and the abstract, introduction, limitations
and conclusion all state that no prototype has been measured.

\rev{18 (PM benchmark)}
\resp Done (Section~X-C, Table~XXI). We sized a first-order analytic surface-PM
machine (12-slot/10-pole, N42SH) to the same mission, envelope, gear, battery,
MOSFETs, iron-loss coefficients, mechanical loss, efficiency boundary and route
model. Its turns follow the same voltage logic as the SRM. Because it matches
the fidelity of the SRM's analytical model, both SRM models are tabulated. The
result is not favorable to the SRM on efficiency:
\begin{itemize}[nosep]
\item the SPM is about 12 points more efficient at the rated point (90.2\%
against 77.9--78.8\%) and far better at launch;
\item it uses 4\% and 20\% less energy on the hilly and aggressive missions;
\item the SRM uses 18\% less on the flat light-assist mission, because the
SPM's load-independent magnet-driven iron loss dominates at light load;
\item the SPM also dissipates about 20~W when coasting without a freewheel
clutch.
\end{itemize}
We report this as it stands and revised the positioning accordingly. The
SRM's case rests on cost, magnet-free supply chain, robustness and light-assist
use, not on peak efficiency.

\rev{19 (quantitative test of the corner-speed rule)}
\resp Done (Section~IX, Fig.~15), as requested.
\begin{itemize}[nosep]
\item \textbf{Mass.} A linear fit of required current against active mass over
the nine designs (3.9-fold mass range) gives a slope of $-0.06$~A/kg
($-0.2$\%/kg) and $R^2=0.004$, so there is no mass dependence.
\item \textbf{Loop factor.} $k_c=0.60\pm0.03$ (range 0.54--0.66). The rule
$P_c/(\bar k_cV)$ reproduces every design to within 10.3\%.
\item \textbf{Voltage.} Redesigning the frame for 24-, 36- and 48-V packs gives
$I_{pk}\propto V^{-1.002}$ with $k_c=0.59$--$0.60$, which tests the voltage
dependence directly.
\end{itemize}

\rev{20 (broad novelty claim)}
\resp Narrowed exactly as suggested: ``To the author's knowledge, the
combination of mission-derived corner-speed turn selection, FE-based current
sizing, and converter, thermal and control evaluation has not been reported
systematically for a 250-W, 36-V geared SRM pedelec drive.''

\rev{21 (placeholder metadata)}
\resp Corrected. The author block and affiliation are complete, and the running
head no longer contains a placeholder volume, issue or date. Reference~[20]
now cites the author.

\rev{22 (repository reproducibility)}
\resp The code reference now cites a tagged release (v2.0-rev2) of the public
repository. The repository contains:
\begin{itemize}[nosep]
\item the source code, FE scripts and all result data (JSON, NPZ);
\item a pinned environment (\texttt{requirements-lock.txt}: Python 3.11.15,
NumPy 2.4.6, SciPy 1.17.1, Matplotlib 3.11.2, scikit-fem 12.0.2, gmsh 4.15.2,
TeX Live 2023);
\item a \texttt{CITATION.cff} file;
\item step-by-step reproduction instructions in the README;
\item 31 automated tests.
\end{itemize}
\emph{[Note to author: a Zenodo DOI can be minted from the GitHub release; add
it to reference~[20] once created.]}

\rev{23 (air-gap qualification)}
\resp Kept. The tolerance allocation remains labeled ``assumed design budget,
not measured''.

\rev{24 (iron-loss sensitivity)}
\resp Done (Table~XII). Scaling the iron-loss model by 0.7 and 1.3, with
operating points re-optimized, moves rated efficiency by about $\pm$2 points
($\eta_{sys}$ 75.9--79.9\%) and route consumption by $-4$\% to $+5$\%. The
paper now identifies iron loss as the second-largest remaining uncertainty.

\rev{25 (acoustic noise)}
\resp Added to the limitations: ``Acoustic noise and radial-force vibration are
outside the present modeling scope, and no claim is made that the design is
acoustically acceptable.''

\rev{28 (title)}
\resp Adopted the reviewer's preferred title: \emph{Mission-Driven Corner-Speed
Design of a Low-Voltage Switched Reluctance Motor Drive for Electric Bicycles}.

\rev{(additional correction made by the author)}
\resp While running the coupled-phase solves, we found that the FE solver
reused the previous solution as its starting point whenever the node count
matched, even when the rotor position, and therefore the mesh, had changed. A
converged Newton solution does not depend on its starting point, so all
previously reported FE results are unaffected. We confirmed this by re-solving
the air-gap study, which reproduces the reported values exactly. However, one
coupled solve diverged. The solver now warm-starts only on the same mesh and
falls back to a damped cold start.

\bigskip
All results can be regenerated from tag v2.0-rev2 of the repository
(\texttt{scripts/run\_fea.py}, \texttt{run\_revision.py},
\texttt{run\_revision2.py}).

\end{document}
